\documentclass[11pt]{article}

% ===================== PACKAGES =====================
\usepackage[a4paper,margin=1in]{geometry}
\usepackage{amsmath,amssymb}
\usepackage{enumitem}
\usepackage{xcolor}

% ===================== SOLUTION COMMAND =====================
\newcommand{\solution}{
    \vspace{0.3cm}
    \noindent\textbf{\textcolor{blue}{Solution:}}
    \vspace{0.3cm}
}

% ===================== TITLE =====================
\title{
CSE 400: Fundamentals of Probability in Computing \\
Homework Assignment-6
}
\date{}

\begin{document}
\maketitle

\begin{enumerate}[label=\textbf{Question-\arabic*:}, leftmargin=*]

% ------------------ Q1 ------------------
\item Let the random vector $X = (X_1, X_2)$ have joint density
\[
f_{X_1,X_2}(x_1,x_2)=
\begin{cases}
c(x_1+x_2), & 0<x_1<1,\; 0<x_2<1 \\
0, & \text{otherwise}
\end{cases}
\]

Find:
\begin{enumerate}[label=(\alph*)]
    \item the value of $c$,
    \item the marginal densities $f_{X_1}(x_1)$ and $f_{X_2}(x_2)$,
    \item $P(X_1>0.3,\; X_2>0.5)$,
    \item whether $X_1$ and $X_2$ are independent.
\end{enumerate}

\solution

% -------- (a) --------
\textbf{(a) Find $c$}

\[
\int_0^1 \int_0^1 c(x_1+x_2)\,dx_2\,dx_1 = 1
\]

\[
\Rightarrow c(1)=1 \;\Rightarrow\; \boxed{c=1}
\]

% -------- (b) --------
\textbf{(b) Marginal densities}

\[
f_{X_1}(x_1)=\int_0^1 (x_1+x_2)\,dx_2
= x_1 + \frac{1}{2}
\]

\[
f_{X_2}(x_2)=\int_0^1 (x_1+x_2)\,dx_1
= x_2 + \frac{1}{2}
\]

% -------- (c) --------
\textbf{(c) Probability}

\[
P=\int_{0.3}^{1}\int_{0.5}^{1}(x_1+x_2)\,dx_2\,dx_1
= \boxed{0.49}
\]

% -------- (d) --------
\textbf{(d) Independence}

\[
f_{X_1}f_{X_2} \neq f_{joint}
\Rightarrow \boxed{\text{Not independent}}

\newpage

% ------------------ Q2 ------------------
\item Suppose the random vector $X = (X_1, X_2, X_3)$ is uniformly distributed over the rectangular box $-1<x_1<1,\; 0<x_2<2,\; 1<x_3<3$. Thus,
\[
f_{X_1,X_2,X_3}(x_1,x_2,x_3)=
\begin{cases}
k, & -1<x_1<1,\; 0<x_2<2,\; 1<x_3<3 \\
0, & \text{otherwise}
\end{cases}
\]

Find:
\begin{enumerate}[label=(\alph*)]
    \item the constant $k$,
    \item the joint marginal density $f_{X_1,X_3}(x_1,x_3)$,
    \item the marginal density $f_{X_2}(x_2)$,
    \item the conditional density $f_{X_1|X_2,X_3}(x_1\mid x_2,x_3)$,
    \item whether $X_1, X_2, X_3$ are mutually independent.
\end{enumerate}

\solution

% -------- (a) --------
\textbf{(a) Find $k$}

Volume of box $= 2 \times 2 \times 2 = 8$

\[
k \cdot 8 = 1 \Rightarrow \boxed{k = \frac{1}{8}}
\]

% -------- (b) --------
\textbf{(b) Joint marginal $f_{X_1,X_3}(x_1,x_3)$}

\[
f_{X_1,X_3}(x_1,x_3)=\int_0^2 \frac{1}{8}\,dx_2
= \frac{1}{8}\cdot 2
= \frac{1}{4}
\]

% -------- (c) --------
\textbf{(c) Marginal $f_{X_2}(x_2)$}

\[
f_{X_2}(x_2)=\int_{-1}^{1}\int_{1}^{3} \frac{1}{8}\,dx_3\,dx_1
= \frac{1}{8}\cdot 2 \cdot 2
= \frac{1}{2}
\]

% -------- (d) --------
\textbf{(d) Conditional density}

\[
f_{X_1|X_2,X_3}=\frac{1/8}{1/4}=\frac{1}{2}
\]

% -------- (e) --------
\textbf{(e) Independence}

\[
f_{X_1}=\frac{1}{2},\quad f_{X_2}=\frac{1}{2},\quad f_{X_3}=\frac{1}{2}
\]

\[
f_{X_1}f_{X_2}f_{X_3}=\frac{1}{8}=f_{joint}
\]

\[
\Rightarrow \boxed{\text{Independent}}

\newpage

% ------------------ Q3 ------------------
\item Let $Y = AX$, where
\[
\mu_X =
\begin{bmatrix}1\\-1\\2\end{bmatrix},\quad
A =
\begin{bmatrix}1&2&0\\0&-1&1\\2&0&1\end{bmatrix},\quad
\Sigma_X =
\begin{bmatrix}4&1&0\\1&3&-1\\0&-1&5\end{bmatrix}.
\]

Compute:
\begin{enumerate}[label=(\alph*)]
    \item the mean vector $\mu_Y$,
    \item the covariance matrix $\Sigma_Y$,
    \item the variance of $Y_1+Y_2$.
\end{enumerate}

\solution

\textbf{(a) Mean vector}

\[
\mu_Y = A\mu_X =
\begin{bmatrix}
1 & 2 & 0\\
0 & -1 & 1\\
2 & 0 & 1
\end{bmatrix}
\begin{bmatrix}
1\\
-1\\
2
\end{bmatrix}
=
\begin{bmatrix}
-1\\
3\\
4
\end{bmatrix}
\]

\textbf{(b) Covariance matrix}

\[
\Sigma_Y = A\Sigma_X A^T
\]

\[
\Sigma_Y =
\begin{bmatrix}
20 & -9 & 10\\
-9 & 10 & 4\\
10 & 4 & 21
\end{bmatrix}
\]

\textbf{(c) Variance of $Y_1+Y_2$}

\[
\mathrm{Var}(Y_1+Y_2)
= \Sigma_{11} + \Sigma_{22} + 2\Sigma_{12}
= 20 + 10 + 2(-9)
= 12
\]


\newpage

\item Three sensor readings are modeled by the random vector $X=(X_1,X_2,X_3)^T$ with
\[
\mu_X =
\begin{bmatrix}2\\0\\-1\end{bmatrix},\quad
\Sigma_X =
\begin{bmatrix}6&2&1\\2&5&-1\\1&-1&4\end{bmatrix},\quad
A =
\begin{bmatrix}1&1&0\\0&2&-1\\1&0&1\end{bmatrix}.
\]

A processing unit forms $Y = AX$. Find:
\begin{enumerate}[label=(\alph*)]
    \item $\mu_Y$,
    \item $\Sigma_Y$,
    \item which component of $Y$ has the largest variance.
\end{enumerate}

\solution

\textbf{(a) Mean vector}

\[
\mu_Y = A\mu_X =
\begin{bmatrix}
1 & 1 & 0\\
0 & 2 & -1\\
1 & 0 & 1
\end{bmatrix}
\begin{bmatrix}
2\\
0\\
-1
\end{bmatrix}
=
\begin{bmatrix}
2\\
1\\
1
\end{bmatrix}
\]

\textbf{(b) Covariance matrix}

\[
\Sigma_Y = A\Sigma_X A^T
\]

\[
\Sigma_Y =
\begin{bmatrix}
15 & 11 & 6\\
11 & 28 & -1\\
6 & -1 & 12
\end{bmatrix}
\]

\textbf{(c) Largest variance}

\[
\mathrm{Var}(Y_1)=15,\;
\mathrm{Var}(Y_2)=28,\;
\mathrm{Var}(Y_3)=12
\]

\[
\Rightarrow Y_2 \text{ has the largest variance}
\]

\newpage
% ------------------ Q5 ------------------
\item A random vector $X = (X_1, X_2, X_3)^T$ has mean vector
\[
\mu_X =
\begin{bmatrix}1\\2\\-1\end{bmatrix}
\]
and second-moment matrix $R_{XX} = E[XX^T]$:
\[
R_{XX} =
\begin{bmatrix}5&3&-1\\3&10&-2\\-1&-2&6\end{bmatrix}.
\]

Find the covariance matrix $C_{XX}$.

\solution

formula:
\[
C_{XX} = R_{XX} - \mu_X\mu_X^T
\]

so $\mu_X\mu_X^T$:

\[
\mu_X\mu_X^T =
\begin{bmatrix}1\\2\\-1\end{bmatrix}
\begin{bmatrix}1&2&-1\end{bmatrix}
=
\begin{bmatrix}1&2&-1\\2&4&-2\\-1&-2&1\end{bmatrix}
\]

Subtract:

\[
C_{XX} =
\begin{bmatrix}5&3&-1\\3&10&-2\\-1&-2&6\end{bmatrix}
-
\begin{bmatrix}1&2&-1\\2&4&-2\\-1&-2&1\end{bmatrix}
=
\begin{bmatrix}4&1&0\\1&6&0\\0&0&5\end{bmatrix}
\]

\newpage

% ------------------ Q6 ------------------
\item Suppose a random vector $X = (X_1, X_2, X_3)^T$ has covariance matrix
\[
C_{XX} =
\begin{bmatrix}3&1&-1\\1&4&2\\-1&2&5\end{bmatrix}
\]
and mean vector
\[
\mu_X =
\begin{bmatrix}2\\-1\\1\end{bmatrix}.
\]

Find the correlation matrix $R_{XX} = E[XX^T]$.

\solution

so formula:
\[
R_{XX} = C_{XX} + \mu_X\mu_X^T
\]

Compute $\mu_X\mu_X^T$:

\[
\mu_X\mu_X^T =
\begin{bmatrix}2\\-1\\1\end{bmatrix}
\begin{bmatrix}2&-1&1\end{bmatrix}
=
\begin{bmatrix}4&-2&2\\-2&1&-1\\2&-1&1\end{bmatrix}
\]

Add:

\[
R_{XX} =
\begin{bmatrix}3&1&-1\\1&4&2\\-1&2&5\end{bmatrix}
+
\begin{bmatrix}4&-2&2\\-2&1&-1\\2&-1&1\end{bmatrix}
=
\begin{bmatrix}7&-1&1\\-1&5&1\\1&1&6\end{bmatrix}
\]

\newpage

% ------------------ Q7 ------------------
\item Suppose $(X_1, X_2, X_3)$ is a jointly Gaussian random vector with mean vector
\[
\mu =
\begin{bmatrix}\mu_1\\\mu_2\\\mu_3\end{bmatrix}
\]
and covariance matrix
\[
\Sigma =
\begin{bmatrix}\sigma_1^2&0&0\\0&\sigma_2^2&0\\0&0&\sigma_3^2\end{bmatrix}.
\]

Prove that the joint density factors as $f_{X_1,X_2,X_3}(x_1,x_2,x_3) = f_{X_1}(x_1)\,f_{X_2}(x_2)\,f_{X_3}(x_3)$, and hence conclude that $X_1, X_2, X_3$ are independent.

\solution

Since the covariance matrix is diagonal, all off-diagonal terms are zero.

\[
\text{Cov}(X_i, X_j) = 0 \quad \text{for } i \neq j
\]

For a jointly Gaussian random vector, zero covariance implies independence.

The joint PDF of a multivariate Gaussian is:

\[
f(x_1,x_2,x_3)
=
\frac{1}{(2\pi)^{3/2}|\Sigma|^{1/2}}
\exp\left(
-\frac{1}{2}(x-\mu)^T \Sigma^{-1} (x-\mu)
\right)
\]

Since $\Sigma$ is diagonal,

\[
|\Sigma| = \sigma_1^2 \sigma_2^2 \sigma_3^2
\]

and

\[
\Sigma^{-1} =
\begin{bmatrix}
\frac{1}{\sigma_1^2} & 0 & 0\\
0 & \frac{1}{\sigma_2^2} & 0\\
0 & 0 & \frac{1}{\sigma_3^2}
\end{bmatrix}
\]

So the exponent becomes:

\[
-\frac{(x_1-\mu_1)^2}{2\sigma_1^2}
-\frac{(x_2-\mu_2)^2}{2\sigma_2^2}
-\frac{(x_3-\mu_3)^2}{2\sigma_3^2}
\]

Thus,

\[
f(x_1,x_2,x_3)
=
\frac{1}{(2\pi)^{3/2}\sigma_1\sigma_2\sigma_3}
\exp\left(
-\frac{(x_1-\mu_1)^2}{2\sigma_1^2}
\right)
\exp\left(
-\frac{(x_2-\mu_2)^2}{2\sigma_2^2}
\right)
\exp\left(
-\frac{(x_3-\mu_3)^2}{2\sigma_3^2}
\right)
\]

\[
= f_{X_1}(x_1)\; f_{X_2}(x_2)\; f_{X_3}(x_3)
\]

Hence, the joint PDF factors into the product of marginal PDFs.

\[
\Rightarrow \boxed{X_1, X_2, X_3 \text{ are independent}}
\]

\newpage

% ------------------ Q8 ------------------
\item Let $X = (X_1, X_2, X_3, X_4)^T$ be a 4-variate jointly Gaussian random vector with mean $\mu = (\mu_1,\mu_2,\mu_3,\mu_4)^T$ and covariance matrix $\Sigma = \mathrm{diag}(\sigma_1^2,\sigma_2^2,\sigma_3^2,\sigma_4^2)$.

Show that the joint density is
\[
f_X(x) = \prod_{i=1}^{4}
\frac{1}{\sqrt{2\pi\sigma_i^2}}
\exp\!\left(-\frac{(x_i-\mu_i)^2}{2\sigma_i^2}\right).
\]
Explain clearly why zero covariances imply independence in this jointly Gaussian case.

\solution

\solution

 covariance matrix is diagonal so all covariances are 0.

For jointly Gaussian variables, the joint PDF is:

\[
f_X(x) =
\frac{1}{(2\pi)^2 \sigma_1\sigma_2\sigma_3\sigma_4}
\exp\left(
-\frac{1}{2}\sum_{i=1}^{4}\frac{(x_i-\mu_i)^2}{\sigma_i^2}
\right)
\]

so here

\[
= \prod_{i=1}^{4}
\frac{1}{\sqrt{2\pi}\sigma_i}
\exp\left(-\frac{(x_i-\mu_i)^2}{2\sigma_i^2}\right)
\]

\[
= f_{X_1}(x_1)\,f_{X_2}(x_2)\,f_{X_3}(x_3)\,f_{X_4}(x_4)
\]

\[
\Rightarrow \boxed{X_1, X_2, X_3, X_4 \text{ are independent}}
\]

SO, Zero covariance removes all cross terms, so joint PDF splits into product independence.

\newpage

% ------------------ Q9 ------------------
\item Prove that
\[
\mathrm{Var}(X-Y) = \mathrm{Var}(X) + \mathrm{Var}(Y) - 2\,\mathrm{Cov}(X,Y).
\]

\solution

here variance:

\[
\mathrm{Var}(X-Y) = E\!\left[(X-Y)^2\right] - \left(E[X-Y]\right)^2
\]

we will expand $(X-Y)^2 = X^2 - 2XY + Y^2$:

\[
E[(X-Y)^2] = E[X^2] + E[Y^2] - 2E[XY]
\]


\[
\left(E[X-Y]\right)^2 = \left(E[X]-E[Y]\right)^2 = E[X]^2 + E[Y]^2 - 2E[X]E[Y]
\]

Subtract:

\[
\mathrm{Var}(X-Y)
= \bigl(E[X^2]-E[X]^2\bigr) + \bigl(E[Y^2]-E[Y]^2\bigr) - 2\bigl(E[XY]-E[X]E[Y]\bigr)
\]

\[
\boxed{\mathrm{Var}(X-Y) = \mathrm{Var}(X) + \mathrm{Var}(Y) - 2\,\mathrm{Cov}(X,Y)}
\]

\newpage
\item A received signal is modeled as $R = S + N_1 - N_2$, where $S$ is the transmitted signal and $N_1, N_2$ are noise terms.

Using variance and covariance properties, derive an expression for $\mathrm{Var}(R)$ in terms of $\mathrm{Var}(S)$, $\mathrm{Var}(N_1)$, $\mathrm{Var}(N_2)$, and the pairwise covariances. Then simplify when $S, N_1, N_2$ are pairwise uncorrelated.

\solution

Given,
\[
R = S + N_1 - N_2
\]

Using the formula:
\[
\mathrm{Var}(aX + bY + cZ)
= a^2\mathrm{Var}(X) + b^2\mathrm{Var}(Y) + c^2\mathrm{Var}(Z)
+ 2ab\,\mathrm{Cov}(X,Y)
+ 2ac\,\mathrm{Cov}(X,Z)
+ 2bc\,\mathrm{Cov}(Y,Z)
\]

Here,
\[
R = 1\cdot S + 1\cdot N_1 + (-1)\cdot N_2
\]

So,

\[
\mathrm{Var}(R)
= \mathrm{Var}(S) + \mathrm{Var}(N_1) + \mathrm{Var}(N_2)
+ 2\mathrm{Cov}(S,N_1)
- 2\mathrm{Cov}(S,N_2)
- 2\mathrm{Cov}(N_1,N_2)
\]

\textbf{If $S, N_1, N_2$ are pairwise uncorrelated:}

\[
\mathrm{Cov}(S,N_1) = 0,\quad
\mathrm{Cov}(S,N_2) = 0,\quad
\mathrm{Cov}(N_1,N_2) = 0
\]

Hence,

\[
\mathrm{Var}(R)
= \mathrm{Var}(S) + \mathrm{Var}(N_1) + \mathrm{Var}(N_2)
\]

\[
\Rightarrow \boxed{\mathrm{Var}(R)=\mathrm{Var}(S)+\mathrm{Var}(N_1)+\mathrm{Var}(N_2)}
\]
\end{enumerate}

\end{document}